\section{Viewing displacement results}
\label{sec:results}

After a run, the canvas draws the selected field over the image, with a colorbar at its right
edge (Figure~\ref{fig:displacement}). Three groups of the right sidebar decide what is drawn and
how: \ui{FIELD}, \ui{VISUALIZATION} and \ui{PHYSICAL UNITS} (Figure~\ref{fig:display}). The
\ui{LOG} at the bottom of the sidebar records what the application does.

\screenshot[0.42\textwidth]{fig-display.png}{The \ui{FIELD}, \ui{VISUALIZATION} and
\ui{PHYSICAL UNITS} groups of the right sidebar after a run.}{fig:display}

Every field is Lagrangian: a value belongs to a material point and is the change from frame~1 to
the frame on show. Frame~1 is the reference, so its displacement and residual are zero.

\subsection{FIELD}

Four buttons choose the field; one is always selected.

\begin{description}[leftmargin=1.2cm,style=nextline]
  \item[Disp U]
    Horizontal displacement $u$, positive to the right.
  \item[Disp V]
    Vertical displacement $v$, positive downward, as image rows count.
  \item[Magnitude]
    $\sqrt{u^2 + v^2}$, the length of the displacement.
  \item[Residual]
    The grey-level mismatch left after the deformed image is mapped back onto frame~1 with the
    measured displacement: a map of how well each point was tracked.
\end{description}

Displacements are in pixels, or in a physical length unit when \ui{Use physical units} is on
(Section~\ref{sec:physical-units}). The residual is in grey levels.

\screenshot[0.45\textwidth]{fig-displacement.png}{Vertical displacement \ui{Disp V} of a plate
with two holes in tension, drawn on the deformed last frame (\ui{Show on}: \ui{Deformed frame})
with the background image. The colorbar gives the field and its unit, V~(px), and the range of
the colours.}{fig:displacement}

\subsubsection*{The residual}
\label{sec:residual}

For each point $\mathbf{X}$ of the region of interest, RAFTcorr reads the current frame at the
displaced position $\mathbf{X} + \mathbf{u}(\mathbf{X})$ and subtracts the frame-1 image at
$\mathbf{X}$. It removes the local mean of this difference and takes its local root mean square,
both over a Gaussian window of $\sigma = 4$\,px. The result is in grey levels on the 0 to 255
scale to which RAFTcorr reads every image (Section~\ref{sec:images}). It is always computed
against frame~1 with the total displacement, whatever the tracking mode.

\begin{itemize}
  \item Where the displacement maps the current frame back onto frame~1, the residual is small,
        at the level of the image noise.
  \item It grows slowly as the specimen deforms (the speckles change shape) and with changes of
        lighting. Because the local mean is removed, a change of brightness that varies slowly
        across the image does not raise it.
  \item It is large where tracking failed, and where the images themselves differ: a crack that
        opens, a reflection, speckle that flakes off, a surface that leaves the focal plane.
\end{itemize}

Read the residual by comparison, between frames and between areas of one frame, rather than
against a fixed value. A residual that rises over the whole region from one frame to the next
means tracking was lost there; this is what the tracking-quality check detects
(Section~\ref{sec:tracking-quality}). A residual that rises in one area points to that area.
Points without a displacement have no residual. The default colormap of the residual is
\file{magma} (Figure~\ref{fig:residual}).

\screenshot[0.45\textwidth]{fig-residual.png}{The residual of the frame of
Figure~\ref{fig:displacement}, in grey levels.}{fig:residual}

\subsection{Where the field is drawn}

\begin{description}[leftmargin=1.2cm,style=nextline]
  \item[Show on]
    \ui{Deformed frame} (the default) draws each value where its material point is in the
    current frame, over the image of that frame. To do so, RAFTcorr inverts the measured
    displacement: for every pixel of the current frame it finds the point of frame~1 that lands
    there and draws that point's value. Pixels where no point of the region of interest lands
    stay empty. \ui{Reference frame} draws each value at the frame-1 position of its point,
    over the image of frame~1. The values are the same; only their positions differ. On
    frame~1 both views are the same.

    Use \ui{Deformed frame} to check that the field follows the specimen: its edges and holes
    should match the image underneath. Use \ui{Reference frame} to compare frames point by point:
    a point stays at one place on the screen while you play the sequence.
  \item[Show background image]
    On by default. Untick it to see the field alone. \ui{Show on} still applies: the positions
    of the points differ between the two views even with nothing behind them.
  \item[Hidden background]
    What replaces the image while it is hidden: \ui{White} (the default), \ui{Black} or
    \ui{Transparent}, which shows the background of the canvas. The list is enabled only while
    \ui{Show background image} is off. In exported pictures, transparency is kept for PNG and
    TIFF; other formats get white.
  \item[Hide untracked regions]
    On by default. A point of the region of interest is untracked in a frame when it has no
    result there: material that left the field of view, that moved into an area that the region
    of interest of a later frame excludes (a crack, a hole), or that fell outside the area
    measured from a reference frame in incremental tracking. With the option on, such points are left
    transparent. With it off, the \ui{Reference frame} view fills them with the value of the
    nearest tracked point inside the region of interest. On the \ui{Deformed frame} view an
    untracked point has no position to be drawn at, so it stays empty either way.
    The choice is also used for exported images and animations; exported data files always
    keep NaN at untracked points.
\end{description}

\subsection{VISUALIZATION}

\begin{description}[leftmargin=1.2cm,style=nextline]
  \item[Colormap]
    The list offers ten colormaps: \file{jet}, \file{viridis}, \file{turbo}, \file{coolwarm},
    \file{plasma}, \file{inferno}, \file{magma}, \file{RdBu\_r}, \file{seismic} and
    \file{black\_rainbow}. The default is
    \file{turbo}, and \file{magma} for the residual. \file{coolwarm}, \file{RdBu\_r} and
    \file{seismic} are diverging maps: with them, the automatic range is symmetric about zero,
    so white means zero. Choose one of them to see the sign of $u$ or $v$ at a glance.
    \file{black\_rainbow} is \file{jet} with its green band replaced by black. The black sits at
    the middle of the colour range, not at zero, and the automatic range is not symmetric with
    this map: give it a \ui{Fixed} range symmetric about zero if black should mean zero.
  \item[Range]
    \ui{Auto} (the default) sets the colour range to the 2nd and 98th percentiles of the values
    drawn on the current frame, after the unit conversion, and computes it again for every frame.
    The extreme 2\,\% at each end take the end colours. \ui{Fixed} keeps \ui{Min} and \ui{Max}
    for every frame. When you switch to \ui{Fixed}, \ui{Min} and \ui{Max} start from the range
    on screen. The boxes hold three decimals, from $-1000$ to $1000$.
  \item[Opacity]
    The opacity of the field over the image, from 0 (transparent) to 100 (opaque); default 70.
\end{description}

The colormap and the range are remembered for each field: switching from \ui{Disp U} to
\ui{Disp V} restores what you last set for \ui{Disp V}. The strain window shares these settings
for its \ui{U}, \ui{V} and \ui{Magnitude} fields (Section~\ref{sec:strain}).

\begin{tipbox}
In \ui{Auto}, each frame has its own range, so one colour means different values on different
frames. Use \ui{Fixed} when you compare frames, play the sequence, or make figures and
animations.
\end{tipbox}

\subsection{PHYSICAL UNITS}
\label{sec:physical-units}

\begin{description}[leftmargin=1.2cm,style=nextline]
  \item[Use physical units]
    Off by default: displacements are in pixels and velocity in pixels per frame. The two
    fields below are enabled only while it is on.
  \item[Pixel size]
    The physical size of one image pixel, with its unit: nm, \textmu m, mm, cm, m or inch,
    followed by \ui{/ px}. Default 1.0\,mm/px. The box holds four decimals: pick the unit in which
    your pixel size has enough digits (5.12\,\textmu m rather than 0.0051\,mm).
  \item[Frame rate]
    The acquisition rate of the images, in frames per second (\ui{fps}); default 1.000\,fps.
\end{description}

A line below the fields shows the resulting units, for example \emph{Disp: mm Velocity: mm/s},
or \emph{Disp: px Velocity: px/fr} while the option is off. With the option on:

\begin{itemize}
  \item \ui{Disp U}, \ui{Disp V} and \ui{Magnitude} are multiplied by the pixel size and shown
        in its unit, on the colorbar, in the cursor read-out and in \ui{Min} and \ui{Max};
  \item the \ui{Velocity} field of the strain window (Section~\ref{sec:strain-fields}) is
        multiplied by the pixel size and the frame rate and shown in the length unit per second;
  \item the \ui{Analysis} tab can plot against time (Section~\ref{sec:analysis});
  \item the residual stays in grey levels, strains stay dimensionless and angles stay in
        degrees.
\end{itemize}

The setting changes what is shown, not the results: data files keep displacements in pixels and
record the pixel size and the frame rate (Section~\ref{sec:export}). The main window and the
strain window share one setting; a change in either shows in both.

\begin{notebox}
A \ui{Fixed} range keeps its numbers when you switch units. After you turn \ui{Use physical
units} on or off, or change the pixel size, set \ui{Min} and \ui{Max} again or go back to
\ui{Auto}.
\end{notebox}

\subsection{The colorbar}

The colorbar sits at the right edge of the canvas. Its title names the field and its unit
(\emph{U~(px)}, \emph{V~(mm)}, \emph{Residual (grey levels)}), and five ticks run from \ui{Min}
at the bottom to \ui{Max} at the top. It follows the range of the frame on show and disappears
when no field is drawn: on a frame that was not computed (after \ui{Cancel}), and while you edit
a region of interest. The cursor read-out and the frame navigator are described in
Section~\ref{sec:canvas}.

\subsection{LOG}

The \ui{LOG} at the bottom of the right sidebar lists the messages of the application, each with
the time (hours, minutes and seconds). It scrolls to the newest line by itself. Errors are red
and completed steps green. % CHECK: warnings sent with level "warning" (engine, Stopped., OOM retry) show in the info colour, only level "warn" is amber; see report.
A run writes its start, the model, the tiling, the frames with a low coverage, the
tracking-quality note and its end there (Section~\ref{sec:running}); an error writes its full
traceback there. \ui{Clear} empties the \ui{LOG}; it changes nothing else. The strain window has
a \ui{LOG} of its own.
